\documentclass[12pt]{article}
\usepackage[margin=1in]{geometry}
\usepackage{amsmath}
\usepackage{fancyhdr}
\fancypagestyle{plain}{%
  \fancyhf{}
  \fancyhead[L]{POLISCI 450A}
  \fancyhead[R]{Fall 2026}
  %\fancyfoot[C]{\thepage}
}
\usepackage{xcolor}
\usepackage{colortbl}
\usepackage{multirow}
\usepackage{amsmath}
\usepackage{amsfonts}
\usepackage{epsfig}
\usepackage{fancyvrb}
\usepackage{bm}

\renewcommand\r{\right}
\renewcommand\l{\left}
\definecolor{dgreen}{rgb}{0,.6,0.8}
\newtheorem{assumption}{Assumption}
\newtheorem{iass}{Identification Assumption}
\newtheorem{ires}{Identification Result}
\newtheorem{estm}{Estimand}
\newtheorem{esti}{Estimator}
\newtheorem{question}{Question}
\newtheorem{answer}{Answer}
\newtheorem{proposition}{Proposition}
\newcommand{\bs}{\boldsymbol}
\newcommand{\mb}{\mathbf}
\newcommand{\real}{\mathbb{R}}
\renewcommand{\u}{\ensuremath{\mb{u}}}
\renewcommand{\v}{\ensuremath{\mb{v}}}
\newcommand{\U}{\ensuremath{\mb{U}}}
\newcommand{\M}{\ensuremath{\mb{M}}}
\newcommand{\X}{\ensuremath{\mb{X}}}
\newcommand{\x}{\ensuremath{\mb{x}}}
\newcommand{\y}{\ensuremath{\mb{y}}}
\renewcommand{\b}{\ensuremath{\bs{\beta}}}
\newcommand{\e}{\ensuremath{\bs{\epsilon}}}
\newcommand{\bhat}{\ensuremath{\hat{\bs{\beta}}}}
\newcommand{\XX}{\ensuremath{\X'\X}}
\newcommand{\XXinv}{\ensuremath{\left(\XX\right)^{-1}}}
\newcommand{\hatsig}{\ensuremath{\hat{\sigma}^2}}
\newcommand{\sig}{\ensuremath{\sigma^2}}
\newcommand{\hatmat}{\ensuremath{\X \XXinv \X'}}
\newcommand{\ehat}{\ensuremath{\hat{\e}}}
\newcommand{\tr}{\text{tr}}
\newcommand{\Sig}{\ensuremath{\bs{\Sigma}}}
\newcommand{\SigInv}{\ensuremath{\bs{\Sigma^{-1}}}}
\newcommand{\W}{\ensuremath{\mb{W}}}
\newcommand{\cN}{\mathcal{N}}
\newcommand\spacingset[1]{\renewcommand{\baselinestretch}{#1}\small\normalsize}
\newcommand{\var}{\text{Var}}
\newcommand{\E}{\mathbb{E}}
\newcommand{\cov}{\mbox{Cov}}
\newcommand{\Cov}{\mbox{Cov}}
\newcommand{\cor}{\mbox{Cor}}
\newcommand{\avepop}{\frac{1}{n}\sum_{i=1}^n}
\newcommand{\indep}{{\bot\negthickspace\negthickspace\bot}}
\newcommand\iid{\stackrel{\rm i.i.d.}{\sim}}
\DeclareMathOperator*\plim{plim}


\usepackage{titlesec}
\titleformat{\section}{\large\bfseries}{\thesection}{1em}{}
\usepackage{titling}
\setlength{\droptitle}{-3em}   % move title up
\usepackage{setspace}

\title{\Large Lecture 6-3: Violations of Homoskedasticity}
\author{}
\date{}
\begin{document}
\maketitle
\vspace{-6em}

\noindent \textbf{Goal:} Diagnose and correct violations of homoskedasticity in regression, focusing on heteroskedasticity, clustering, and serial correlation.\medskip

\begin{itemize}\setlength{\itemsep}{1em}
\item Violations of homoskedasticity
\item Heteroskedasticity
\begin{itemize}
\item Definition
\item Consequences for inference
\item Visual diagnostics and residual plots
\item Formal tests (Breusch-Pagan, White)
\item Remedies: transformations, WLS, robust errors
\item White’s heteroskedasticity-robust SE
\end{itemize}
\item Clustering
\begin{itemize}
\item Concept of cluster dependence
\item Error variance matrix structure
\item Estimation of cluster-robust SE
\end{itemize}
\item Serial correlation
\begin{itemize}
\item Time dependence in errors
\item AR(1) error structure
\item Variance and covariance of AR errors
\item Detection: plots, Durbin-Watson
\item Corrections: Newey-West HAC, AR adjustments
\end{itemize}
\item Summary
\begin{itemize}
\item Three key violations: heteroskedasticity, clustering, serial correlation
\item Detection through plots and tests
\item Use of robust or cluster-robust standard errors
\end{itemize}
\end{itemize}



% End outline
\clearpage


\section*{Worksheet}


\begin{enumerate}\itemsep15em

% \item Explain why $\hat{\b}$ remains unbiased and consistent under heteroskedasticity?



\item The classic variance estimator is
\[
\widehat V_{0}
= \hat\sigma^2(\X'\X)^{-1}.
\]
The White estimator is
\[
\widehat V_{HC0}
= (\X'\X)^{-1}\X'\mathrm{diag}(\hat u_1^2,\ldots,\hat u_n^2)\X(\X'\X)^{-1}.
\]
Assume (1) $E[u_i\mid \X]=0$; (2) $\var(u_i\mid\X)=\sigma^2$; (3) $\frac{1}{n}\X'\X\to Q\succ0$; and (4) standard regularity conditions. Show that
\[
\plim (n\,\widehat V_0)=\plim (n\,\widehat V_{HC0})=\sigma^2 Q^{-1}.
\]

\noindent\textbf{Hint:}
\[
\frac{1}{n}\sum_{i=1}^n \x_i\x_i'\hat u_i^2
= \frac{1}{n}\sum_{i=1}^n \x_i\x_i' u_i^2 + o_p(1)
\xrightarrow{p}\sigma^2 E[\x_i\x_i']=\sigma^2 Q.
\]

\clearpage

\item Consider the linear regression model
\[
y_{ij} \ = \ \beta_0 + \beta_1 x_{ij} + \varepsilon_{ij} \ = \
\beta_0 + \beta_1 x_{ij} + v_j + u_{ij},
\]
where units $i$ are nested in clusters $j$ and
\[
v_j \iid N(0,\rho\sigma^2),\qquad
u_{ij} \iid N(0,(1-\rho)\sigma^2),
\]
with $v_j$ and $u_{ij}$ mutually independent and independent of $x_{ij}$. The parameter $\rho\in[0,1]$ is called the within-cluster correlation.

Suppose we have $M=2$ clusters and $N=6$ units:
cluster $1$ contains units $i=1,2,3$ and cluster $2$ contains units $i=4,5,6$.
Stack the errors in the order
\[
\bm{\varepsilon}
=
(\varepsilon_{1,1},\varepsilon_{2,1},\varepsilon_{3,1},
\varepsilon_{4,2},\varepsilon_{5,2},\varepsilon_{6,2})'.
\]

\begin{enumerate}
\item Show that $V[\varepsilon_{ij}] = \sigma^2$ for all $i,j$.

\item For two units $i\neq \ell$ in the \emph{same} cluster $j$, compute $\cov(\varepsilon_{ij},\varepsilon_{\ell j}).$

\item For units in \emph{different} clusters $j\neq k$, compute $\cov(\varepsilon_{ij},\varepsilon_{\ell k}).$

\item Using your answers above, write out the $6\times 6$ covariance matrix $\Sigma = V[\bm{\varepsilon}]$ for the given ordering
$(\varepsilon_{1,1},\varepsilon_{2,1},\varepsilon_{3,1},
\varepsilon_{4,2},\varepsilon_{5,2},\varepsilon_{6,2})'$.

\item Compute the correlation $\cor(\varepsilon_{ij},\varepsilon_{\ell j})$
for $i\neq \ell$ in the same cluster $j$ and verify that it equals $\rho$.
\end{enumerate}

\clearpage


\item The general cluster-robust variance estimator can be written as  
\[
\widehat{V}_{CR}[\hat{\b}] = (\X'\X)^{-1}\left(\sum_{j=1}^M \X_j'\hat{\u}_j\hat{\u}_j'\X_j\right)(\X'\X)^{-1},
\]
where $\hat{\u}_j$ is the vector of residuals in cluster $j$.  
\textbf{Explain} how this generalizes the White estimator.



% \item Show that
% \[
% DW \xrightarrow{p} 2(1 - \hat{\rho}),
% \]
% where $\hat{\rho} = \frac{\sum_{t=2}^{T}\hat{u}_t\hat{u}_{t-1}}{\sum_{t=1}^{T}\hat{u}_t^2}$ is the sample autocorrelation of the residuals (the OLS estimate of $\rho$).  
% In other words, as $T\to\infty$, the Durbin–Watson statistic converges in probability to $2(1-\hat{\rho})$. Explain: (1) if $DW\approx2$ then $\widehat{\rho}\approx0$; (2) if $DW<1$ we have serious positive serial correlation; and (3) if $DW>3$ we have serious negative serial correlation.

\end{enumerate}


\end{document}
